%% LyX 2.5.3 created this file.  For more info, see https://www.lyx.org/.
%% Do not edit unless you really know what you are doing.
\documentclass[12pt,english]{article}
\renewcommand{\rmdefault}{ptm}
\usepackage[T1]{fontenc}
\usepackage[utf8]{inputenc}
\usepackage{amsmath}
\usepackage{graphicx}
\usepackage{cancel}

\makeatletter
%%%%%%%%%%%%%%%%%%%%%%%%%%%%%% Textclass specific LaTeX commands.
\usepackage[natbibapa]{apacite}

\@ifundefined{date}{}{\date{}}
\makeatother

\usepackage{babel}
\begin{document}
\title{The role of demographic feedback on the evolution of developmental
windows}
\maketitle

\part{Main Manuscript}

\section{Introduction}

There is an ongoing mass extinction event and, as many of the relevant
anthropogenic stressors continue to worsen---not least but not limited
to rising global temperatures and the consequent risks of climate
system breakdown (CITE)---biological models have an important job
to do in predicting the demographic and adaptive responses of species
and ecosystems to these changing conditions. There remains much we
can do, in principle, to improve these models and so generate more
robust and informative guidance for policy makers. For instance, models
of demographic change have yet to synthesise reports from the evolutionary
theory literature of the great speeds (within ecological timescales)
at which adaptation can occur and of the important role of plasticity
and agency of individuals in directing this adaptation (CITE). Benton
(2012) (CITE), shows how the complexities of real populations resulting
from individual-level genetic variation and plasticity can play a
strong role in population responses. We need more models like this,
elucidating what variables are important in population dynamics, in
order to develop appropriate theory, generate accurate empirical predictions,
and present robust and effective policy guidance (CITE). Research
within behavioural ecology has increasingly focused on the roles of
uncertainty, state-dependent decision making, and learning; with implications
for the evolution of various behavioural and developmental strategies
(CITE CITE CITE), but their roles in population dynamics and the implications
for us remain largely speculative (CITE). 

\subsection{Main question with examples}

The main question is, given the predictability and severity of food
limitations in an organism's environment, over what length of time
do individuals invest in developing higher reproductive quality at
an opportunity cost to investing in growth to exploit periods of abundant
food?

Developmental windows are periods in development where sensitivity
to cues is heightened.

Walasek et al. (2021), Walasek et al. (2022) find that the timing
of when sensitive periods emerge depends on the reliability of cues
increases across ontogeny, and their results suggest that both information
and phenotypic states are modeled, with incremental development in
fluctuating environments favouring residual development across ontogeny.

Further, Walasek et al (2024) review current theory and empirical
literature and suggest future models could better integrate with empirical
research by incorporating phenomena found in empirical systems into
our models, and design more models that can be tailored to the life
histories of specific species. They encourage for instance that models
that incorporate the life history trade-offs we see in the development
of alternative phenotypes in bulb mites (between scramblers and fighters)
(Smallegange, 2023; Smallegange et al., 2019), and it happens that
authors interested in these taxa are also interested in bridging the
gap in research between demographic modelling and agency with respect
to development shaping demography (Smallegange ).

\subsection{Developmental windows and sensitive periods}

\subsection{Demography, agency and life-history}

Feedback between plastic responses and population dynamics has been
explored for fixed reaction norms by Brass et al. (2021) who showed
that \_\_\_. However, it is unclear how these eco-devo dynamics will
play out over evolutionary time.

\subsection{Hypothesis}

\subsection{Methods}

\section{The Model Outline}

We model a large population of asexual individuals, distributed in
a metapopulation across habitat patches of equal size and quality,
but each with a population density state with transition probabilities
determined by the local population's growth rate. Each individual
has three state dimensions---size ($s$), patch information ($\pi$),
and reproductive quality ($m$)---upon which their developmental
strategy is dependent, $x(s,\pi,m)$, which determines the fraction
of ingested energy allocated to developing greater reproductive quality
over growth to increase their size state (and so ultimately the local
population growth rate). Size determines whether or not an individual
is mature given they can only reproduce when they reach their maximum
size $S$. An individual's patch information state gives them the
subjective probabilities regarding each possible population density
state of their habitat patch. Finally, an individual's reproductive
quality state $m$ determines their rate of reproduction at maturity. 

We assume mutations are rare and of small effect such that effectively
all individuals across the metapopulation have the same state-dependent
developmental strategy (henceforth referred to as the resident strategy).
Given the resident strategy we then do a forward projection using
a Markov chain until demographic stability is reached at the metapopulation
level (i.e. among patches). Given the probability distribution of
population density states of patches, and assuming individuals have
complete information as to this probability distribution and of the
transition probabilities between patch pop. density states, we find
the best response strategy for an invading individual (i.e. the developmental
strategy that maximises their expected fitness).

\subsection{The Life Cycle and Demographic Feedback}

We model the life cycle of both individuals and of the patches they
reside in (see figure 1).
\begin{enumerate}
\item Newborns disperse into random patches, with accurate prior information
as to the probability distribution over patch population density states,
and causing transitions in patch population density states.
\item Individuals forage, with food availability being a random variable,
normally distributed, with mean determined by the population density.
\item Individuals update their patch information state given informativeness
of food availability.
\item Immature individuals allocate a proportion of their ingested energy
towards increasing their future reproductive quality given their developmental
strategy, the rest of the energy goes towards maturation (increasing
their size). Mature individuals do not grow, investing all ingested
energy towards current reproduction.
\item Individuals die given a baseline probability of death from exogenous
factors and given perturbations in patch quality (e.g. disease outbreaks)
influenced by the patch population density state, and in turn causing
transitions of patch states to lower population densities.
\item Individuals that are of maximum size are mature and reproduce at rate
determined by their reproductive quality state.
\end{enumerate}
\begin{center}
\begin{figure}
\centering{}\includegraphics[width=1\textwidth]{T:/tjcwo/OneDrive - University of Exeter/Documents/PostDocs/Model_windows/Lifecycle.png}\caption{The model structure, including the life cycle of individuals in purple
and the feedback loops acting on patches in green.}
\end{figure}
\par\end{center}

\subsection{State dimensions}
\begin{itemize}
\item Patch population density state, $i\in\{0,1,2\}$, corresponding to
low, medium, and high population densities respectively. We assume
that $i$ determines the probability of patch quality decline and
determining the availability of food.
\item Size state, $s\in\{1,2,...,S\}$, where $S$ is their size at maturity.
\item Patch belief state, $\pi=\{\pi_{0},\pi_{1},\pi_{2}\}$, is an individual's
implicit and subjective probability distribution over patch population
density states, where $\pi_{i}=P(i)$, $\pi_{i}\in\{0/L,1/L,...,L/L\}$,
and $\pi_{0}+\pi_{1}+\pi_{2}=1$, where $L$ determines the resolution
of the belief state. Thus, note that $\pi$ is best represented as
co-ordinates within a discretised probability 2-simplex (i.e. an equilateral
triangle where each vertex represents each possible state of certainty),
rather than representing $\pi$ with a 3D cube (which would include
impossible probability distributions and so waste state space).
\item Reproductive quality state, $m\in\{0,1/M,2/M,...,M/M\}$, corresponding
to the quality of an individual's physiology and so the mass of their
offspring. Further, any investment in quality above the quality threshold
$m_{switch}$, results in an individual developing an alternative
fighter phenotype, giving additional reproductive benefit $r_{fighter}$.
\end{itemize}

\part{Appendix}

\section{Finding Demographic Stability}

Here we describe pseudocode for calculating the demographically stable
probability distribution of population densities across patches and
of the transition probabilities between patch density states. Patches
have a population size state $i\in\{0,1\}$. The actual population
size for any given state $q_{i}$ is defined as
\begin{equation}
q_{i}=n_{low}+n_{max}\cdot i\label{eq:}
\end{equation}

where $n_{max}$ is the maximum population size of any patch. 

\subsection{Reproduction}

Adults, individuals of size state $s=S$, reproduce asexually at a
rate dependent on their reproductive quality state. We here define
$R(s,m)$, the expected reproductive rate of an individual of state
$\{s,m\}$ in a patch of density $q_{i}$, as 
\begin{equation}
R(s,m,i)=\begin{cases}
\text{if \ensuremath{s\geq S} and }m\geq m_{switch}: & \sum_{g}p(g|q_{i})\cdot g\cdot(r_{b}+m\cdot r_{m}+r_{fighter})\\
\text{if \ensuremath{s\geq S} and }m<m_{switch} & \sum_{g}p(g|q_{i})\cdot g\cdot(r_{b}+m\cdot r_{m})\\
\text{else}: & 0
\end{cases}\label{eq:-1-1}
\end{equation}
where S is their size at maturity, $m_{switch}$ is the threshold
investment in quality above which the fighter polyphenism is expressed,
$r_{b}$ is the baseline reproductive rate, $r_{m}$ is the reproductive
benefit of greater somatic quality, and $r_{fighter}$ is reproductive
benefit of expressing the fighter polymorphism allowing them to extract
auxiliary resources for reproduction. To make solving the model tractable
we must assume the fighter phenotype does not influence competition
for the shared limited resource represented by $g$.

\subsection{Dispersal}

\subsubsection{Migrant Physiology}

Ideally, individuals would be born with precise and accurate prior
information regarding the probability distribution over the population
density states of patches. Individuals would then update their information
state given cues for the population density of their patch which come
from the availability of food, see section 1.2. Individuals are born
at size state $0$ and reproductive quality state 0, and a belief
state over population densities, $\forall i$ of
\begin{equation}
\pi_{i}\approx h(q_{i})\label{eq:-3}
\end{equation}
However, to form a discrete grid of states that we can solve a Markov
decision process for we need interpolate each $\pi_{i}$ onto a grid
of possible discrete probability values, $\pi_{i}\in\{0/L,1/L,2/L,...,L/L\}$,
where $L$ is an integer determining the fitness of the grid.then
needs to be interpolated onto a discrete grid of possible values of
$\pi_{i}$ which (because it represent a probability distribution)
lies on a unit simplex (i.e. all values are positive and they add
up to 1), requiring a grid of barycentric co-ordinates. To this end,
the first step is to decompose each $\pi^{\prime}_{i}$ into it's
integer and fractional parts, with the integer part defined by finding
the floor integer after scaling by the maximum value $L$ 
\[
n_{i}=\lfloor h(q_{i})\cdot L\rfloor,
\]
and subsequently the fractional part defined as
\[
f_{i}=h(q_{i})\cdot L-\iota_{i}.
\]

Note that just as 
\[
\sum_{i}h(q_{i})=1
\]
so does
\[
L=\sum_{i}h(q_{i})\cdot L,
\]
and so 
\[
L=\sum_{i}n_{i}+f_{i}
\]
and therefore the sum of the fractional parts must equal an integer
because $L$ is an integer. Therefore, assuming $\pi=\{\pi_{0},\pi_{1},\pi_{2}\}$,
the fractional parts must add up to 0, 1 or 2. And so for all states
the proportion of dispersers that are of state $\{s,\pi,m\}$ we first
set to zero, $\tilde{\rho}(s,\pi,m)=0$. Then we update given these
fractional parts. If the fractional parts add up to $0$, all offspring
have initial state
\[
\tilde{\rho}(s=1,\pi=\{n_{1},n_{2},n_{3}\},m)=1
\]
Else, if the fractional parts add up to 1 we know exactly one dimension
of their prior belief state has to be rounded up, with weightings
given to their respective fractional parts and so 
\[
\tilde{\rho}(s=1,\pi=\{n_{0}+1,n_{1},n_{2}\},m=0)=f_{0}
\]
\[
\tilde{\rho}(s=1,\pi=\{n_{0},n_{1}+1,n_{2}\},m=0)=f_{1},
\]
\[
\tilde{\rho}(s=1,\pi=\{n_{0},n_{1},n_{2}+1\},m=0)=f_{2}.
\]
Else and finally, if the fractional parts add up to 2, exactly 2 dimensions
have to be rounded up
\[
\tilde{\rho}(s=1,\pi=\{n_{0},n_{1}+1,n_{2}+1\},m=0)=1-f_{0},
\]
\[
\tilde{\rho}(s=1,\pi=\{n_{0}+1,n_{1},n_{2}+1\},m=0)=1-f_{1},
\]
\[
\tilde{\rho}(s=1,\pi=\{n_{0}+1,n_{1}+1,n_{2}\},m=0)=1-f_{2}.
\]


\subsubsection{Patch Densities transition probabilities}

Assuming a dispersal rate of $\mu$ and a fixed number of patches
$N$, the number of individuals arriving in each patch of population
density $q_{i}$ is
\begin{equation}
\alpha(q_{i})=(1-\mu)\cdot q_{i}\cdot r(q_{i})+\frac{\mu\cdot\cancel{N}\cdot\sum^{I}_{j=0}h(q_{j})\cdot q_{j}\cdot r(q_{j})}{\cancel{N}}\label{eq:-4}
\end{equation}
where $r(q)$ is the expected reproductive rate of an individual in
a patch of population density $q_{i}$ and where $h(q_{i})$ is the
proportion of patches that are of density $q_{i}$. Note that $N$
in the numerator of the fraction is there to scale the number of migrants
given the number of patches producing them, but that $N$ gets canceled
out as these migrants are dispersed across these patches (see the
denominator), telling us that the number of immigrants to a patch
is not dependent on the number of patches; only the frequency distribution
over patch states matters. The expected reproductive rate of an individual
in a patch of population density $q_{i}$ is defined as 
\begin{equation}
r(q_{i})=\sum^{S}_{s}\sum^{\Pi}_{\pi}\sum^{M}_{m}[\rho_{i}(s,\pi,m)\cdot R(s,m,i)]\label{eq:-5}
\end{equation}
where $\rho_{i}(s,\pi,m)$ is the probability that an individual will
be of state $\{s,\pi,m\}$ given that they are in a patch of population
density $q_{i}$, and $R(s,m,i)$ is the reproductive rate of an individual
of state $\{s,m\}$ in a patch of density $q_{i}$ (see section 1.4). 

Following arrivals, the new population density of patches of prior
density $q_{i}$ become 
\begin{equation}
q^{\prime}_{i}=min(q_{I},q_{i}+\alpha(q_{i}))\label{eq:-6}
\end{equation}
We must then interpolate the posterior population density $q^{\prime}_{i}$
back on to the grid of density states, with the nearest positions
being
\begin{equation}
d_{1}(i)=\lfloor\frac{q^{\prime}_{i}}{n_{max}}\cdot I\rfloor\label{eq:-7}
\end{equation}
and 
\begin{equation}
d_{2}(i)=d_{1}(i)+1.\label{eq:-8}
\end{equation}
Therefore patches of population density $q_{i}$ transition to patches
of density state $q_{d_{2}(i)}$ with probability 
\begin{equation}
p_{2}(i)=(\frac{q^{\prime}_{i}}{n_{max}}\cdot I)-d_{1}(i),\label{eq:-9}
\end{equation}
and transition to density state $q_{d_{1}(i)}$ with probability 
\begin{equation}
p_{1}(i)=1-p_{2}(i).\label{eq:-10}
\end{equation}

This then returns the transition probabilities between patch density
states given dispersal $p_{births}(q_{d_{1}(i)}|q_{i})$, which at
demographic stability will be used for finding the best response strategy.
For all $q_{i}$ $p_{births}(q_{d_{1}(i)}|q_{i})=0$ and then
\begin{itemize}
\item Add $p_{1}(i)$ to $p_{births}(q_{d_{1}(i)}|q_{i})$
\item Add $p_{2}(i)$ to $p_{births}(q_{d_{2}(i)}|q_{i})$.
\end{itemize}

\subsubsection{Projection}

The posterior frequency distribution over patch population density
states is then defined as 
\begin{equation}
h^{\prime}(q_{i})=\sum_{j}h(q_{j})\cdot p_{births}(q_{i}|q_{j})\label{eq:-11}
\end{equation}

Regarding the distribution of physiological states within patches,
individuals in any patch state $q_{i}$ will have a physiological
state frequency distribution defined given the transition probabilities
between patches and weighted by the change in population size for
these transitions. Thus, $\forall i\forall\{s,\pi,m\}$
\begin{equation}
\rho^{\prime}_{i}(s,\pi,m)=\sum_{j}[h(q_{j})\cdot p_{births}(q_{i}|q_{j})\cdot(\frac{j}{i}\cdot\rho_{j}(s,\pi,m)+\frac{i-j}{i}\cdot\tilde{\rho}(s,\pi,m))]\label{eq:-12}
\end{equation}

where if patch were to transition from a density of $j$ to a density
of $i$, $(\frac{j}{i}\cdot\rho_{i}(s,\pi,m)+\frac{i-j}{i}\cdot\tilde{\rho}(s,\pi,m))$
is the proportion of individuals who are of state $\{s,\pi,m\}$ within
that patch. This does however indicate that different patches will
have different physiological distributions amongst their populations
depending on the historical density states of that patch which would
violate our assumptions of this being a Markov process, so (to maintain
tractability) we assume that individuals have equal access to food
regardless of physiological state.

\subsection{Foraging}

Individuals observe cues for population density in the form of food
availability. Food availability is a random variable influenced by
the population density, all patches produce the same amount of food
which must be shared, so for any given patch of population density
$q_{i}$ individuals receive a food mass of 
\begin{equation}
g\sim\text{Normal}(\frac{\bar{g}}{q_{i}},\sigma_{g})\label{eq:-14}
\end{equation}
where parameter $\bar{g}$ is the mean food mass production of any
patch and parameter $\sigma_{g}$ is the food mass standard deviation.
However, as we are working with discrete mathematics, we must approximate
this normal distribution using a pseudo-normal probability distribution
over a grid of possible food mass values $\boldsymbol{G}=\{0,1,...,G_{max}\}$,
with probabilities for each realised food mass being 
\begin{equation}
\begin{array}{ccc}
p(g|q_{i})=\frac{1}{K}e^{-\frac{(g-\frac{\bar{g}}{q_{i}})^{2}}{2\sigma^{2}_{g}}} & \text{for} & g\in\boldsymbol{G},\end{array}\label{eq:-2}
\end{equation}
where $N$ is a suitable integer that determines the fineness of the
grid of values, and $K$ is 
\begin{equation}
K=\underset{-G}{\overset{G}{\boldsymbol{\sum}}}e^{-\frac{(g-\frac{\bar{g}}{q_{i}})^{2}}{2\sigma^{2}_{g}}}.\label{eq:-15}
\end{equation}


\subsubsection{Observation}

Given the food mass $g$, individuals observe the resulting nutrient
data and update their population density information states (i.e.
their belief states). Assuming Bayesian optimal updating, their posterior
belief is defined $\forall i$ as
\begin{equation}
\pi^{\prime}_{i}=\frac{\pi_{i}\cdot p(g|q_{i})}{\sum_{j}\pi_{j}\cdot p(g|q_{j})}.\label{eq:-16}
\end{equation}
The posterior belief state $\pi^{\prime}_{i}$ then needs to be interpolated
onto a discrete grid of possible values of $\pi_{i}$ which (because
it represent a probability distribution) lies on a unit simplex (i.e.
all values are positive and they add up to 1), requiring a grid of
barycentric co-ordinates. To this end, the first step is to decompose
each $\pi^{\prime}_{i}$ into it's integer and fractional parts, with
the integer part defined by
\[
n_{i}=\lfloor\pi^{\prime}_{i}\cdot L\rfloor,
\]
and subsequently the fractional part defined as
\[
f_{i}=\pi^{\prime}_{i}\cdot L-\iota_{i}.
\]

Note that just as 
\[
\sum_{i}\pi^{\prime}_{i}=1
\]
so does
\[
L=\sum_{i}\pi^{\prime}_{i}\cdot L,
\]
and so 
\[
L=\sum_{i}n_{i}+f_{i}
\]
and therefore the sum of the fractional parts must equal an integer
because $L$ is an integer. Therefore, assuming $\pi=\{\pi_{0},\pi_{1},\pi_{2}\}$,
the fractional parts must add up to 0, 1 or 2. If they add up to $0$
\[
p(\pi^{\prime}=\{n_{1},n_{2},n_{3}\}|\pi,g\}=1
\]
If the fractional parts add up to 1 we know exactly one dimension
of the posterior belief state has to be rounded up, with weightings
given to their respective fractional parts and so 
\[
p(\pi^{\prime}=\{n_{0}+1,n_{1},n_{2}\}|\pi,g)=f_{0}
\]
\[
p(\pi^{\prime}=\{n_{0},n_{1}+1,n_{2}\}|\pi,g)=f_{1},
\]
\[
p(\pi^{\prime}=\{n_{0},n_{1},n_{2}+1\}|\pi,g)=f_{2}.
\]
Else if the fractional parts add up to 2, exactly 2 dimensions have
to be rounded up
\[
p(\pi^{\prime}=\{n_{0},n_{1}+1,n_{2}+1\}|\pi,g)=1-f_{0},
\]
\[
p(\pi^{\prime}=\{n_{0}+1,n_{1},n_{2}+1\}|\pi,g)=1-f_{1},
\]
\[
p(\pi^{\prime}=\{n_{0}+1,n_{1}+1,n_{2}\}|\pi,g)=1-f_{2}.
\]
Finally, for all other possible posterior co-ordinates on the belief
grid, the transition probability is 
\[
p(\pi^{\prime}|\pi,g)=0.
\]


\subsubsection{Development}

Each individual, choose a developmental action given their state $\{s,\pi,m\}$,
given the resident developmental strategy for that state $x(s,\pi^{\prime},m)\in(0,G)$.
Their developmental action determines the threshold mass above which
energy foraged goes to the development of phenotypes of greater reproductive
quality (e.g. fighter morphology), generating a transition in their
size state $s$ and their reproductive quality state $m$. Of the
portion of energy allocated to quality, all energy below a size-dependent
investment threshold goes to maintenance costs. Therefore, if investment
in quality falls below this threshold, given their state-dependent
strategy and their food mass, there is a discounted probability of
survival
\begin{equation}
P(survival|g,s,\pi^{\prime},m)=min(1,\frac{g-\Delta m-\Delta s}{y(s,m)})\label{eq:-44}
\end{equation}
such that if the individual has no energy to allocate to maintenance
they die, where $y(s)$ is the size dependent maintenance cost 
\begin{equation}
y(s,m)=\frac{s^{3}\cdot m}{S^{3}\cdot M}\cdot y_{max}\label{eq:-22}
\end{equation}
where we assume maintenance costs scale by the volume of the individual
and by the quality of their somatic cells and developmental systems
(given their quality state $m$), and where $y_{max}$ is the maximum
maintenance investment. 

The resulting transition probabilities, $p(s^{\prime},\pi^{\prime},m^{\prime}|s,\pi,m)$,
between state $\{s^{\prime},\pi^{\prime},m^{\prime}\}$ and state
$\{s,\pi,m\}$ given the resident developmental strategy $x(s,\pi,m)$
and discounted survival probabilities, are found be doing the following: 
\begin{enumerate}
\item For all $\{s,\pi,m,g\}$
\[
p(s^{\prime},m^{\prime}|s,\pi^{\prime},m,g)=0.
\]
\item For all $g$ and for each state, calculate the changes in quality
and size
\begin{equation}
\Delta m=max(g-x(s,\pi^{\prime},m),0)\label{eq:-20-1}
\end{equation}
where $\Delta m$ is the change in quality state of individuals after
ingesting a food mass of $g$, and where
\begin{equation}
\Delta s=max(g-\Delta m-y(s,m),0),\label{eq:-19-1}
\end{equation}
is the change in size state of individuals after ingesting a food
mass of $g$, given that all energy allocated to growth below the
fixed and state-dependent maintenance cost threshold $y(s,m)$, below
which all food mass is invested in maintenance.
\item These changes then need to be interpolated onto the size and quality
grids, and for all interpolated values we find the resulting discretised
transitions and their respective probabilities
\[
p(s+\Delta s,m+\Delta m|s,\pi^{\prime},m,g)=p(\Delta s)\cdot p(\Delta m)
\]
\end{enumerate}

\subsubsection{Updating the States Distributions}

Given the baseline rate of death $\lambda$ and deaths caused by lack
of investment in maintenance, the posterior density of each state
for patches of each state $i$ is 
\begin{equation}
\hat{\rho^{\prime}_{i}}(s^{\prime},\pi^{\prime},m^{\prime})=\sum_{g}\sum_{s}\sum_{\pi}\sum_{m}p(s^{\prime},m^{\prime}|s,\pi^{\prime},m,g)\cdot P(survival|g,s,\pi^{\prime},m)\cdot p(\pi^{\prime}|\pi,g)\cdot p(g|q_{i})\cdot\rho^{\prime}_{i}(s,\pi,m)\label{eq:-17}
\end{equation}

Where, as a reminder, $\rho^{\prime}_{i}(s,\pi,m)$ is the probability
distribution over states after dispersal within patches of population
size $q_{i}$, $p(g|q_{i})$ is the probability of a food mass of
$g$ in patches of size $q_{i}$, $p(\pi^{\prime}|\pi,g)$ is the
probability of transitioning to posterior information state $\pi^{\prime}$
given this food mass, $P(survival|s,\pi^{\prime},m,g)$ is the probability
of survival of maintenance costs given their developmental strategy
dependent on posterior information state $\pi^{\prime}$, and $p(s^{\prime},m^{\prime}|s,\pi^{\prime},m,g)$
is the probability of transitioning to size state $s^{\prime}$ and
quality state $m^{\prime}$ given the outcomes of development and
given that they have survived maintenance costs.

Given the baseline rate of death $\lambda$ and deaths caused by lack
of investment in maintenance, patches can transition to lower population
density states. Therefore, given the posterior densities of individuals
within patches, the posterior of population size of a patch of density
state $q_{i}$ is 
\[
u=q_{i}\cdot(1-\lambda)\cdot\sum_{s}\sum_{\pi}\sum_{m}\hat{\rho^{\prime}_{i}}(s^{\prime},\pi^{\prime},m^{\prime})
\]
 The posterior distribution of patch states is therefore
\begin{equation}
h^{\hat{\prime}}(q_{i})=\sum^{I}_{j}h(q_{j})\cdot p_{\text{fail}}(q_{i}|q_{j})\label{eq:-23}
\end{equation}
where $p_{\text{fail}}(q_{i}|q_{j})$ is the transition probability
of a patch from state $j$ to $i$ given deaths from maintenance failures,
defined by interpolating $u$ so that where 
\begin{equation}
s_{1}(j)=floor(\frac{u}{n_{max}}\cdot I)\label{eq:-26}
\end{equation}
\begin{equation}
s_{2}(j)=s_{1}(j)+1\label{eq:-27}
\end{equation}
\begin{equation}
p_{2}(j)=(\frac{u}{n_{max}}\cdot I)-s_{1}\label{eq:-28}
\end{equation}
\begin{equation}
p_{1}(j)=1-p_{2}(j)\label{eq:-29}
\end{equation}
and so run the following algorithm for each $q_{j}$:
\begin{enumerate}
\item 
\begin{equation}
\forall q_{i},p_{\text{fail}}(q_{i}|q_{j})=0.\label{eq:-21}
\end{equation}
\item 
\begin{equation}
p_{\text{fail}}(s_{1}(q_{j})|q_{j})=p_{1}(j)+p_{\text{fail}}(s_{1}(q_{j})|q_{j}).\label{eq:-24}
\end{equation}
\item 
\begin{equation}
p_{\text{fail}}(s_{2}(q_{j})|q_{j})=p_{2}(j)+p_{\text{fail}}(s_{2}(q_{j})|q_{j}).\label{eq:-25}
\end{equation}
\end{enumerate}
The posterior probability distribution over individual states for
patches of each population density state is therefore 
\begin{equation}
\rho^{\prime\prime}_{i}(s^{\prime},\pi^{\prime},m^{\prime})=\sum_{q_{j}}h(q_{j})\cdot p_{\text{fail}}(q_{i}|q_{j})\cdot\frac{\hat{\rho^{\prime}_{j}}(s^{\prime},\pi^{\prime},m^{\prime})}{\sum_{s}\sum_{\pi}\sum_{m}\hat{\rho^{\prime}_{j}}(s^{\prime},\pi^{\prime},m^{\prime})}\label{eq:-30}
\end{equation}


\subsection{Patch quality perturbations}

Every patch undergoes momentary perturbations in quality, influenced
by patch population density (e.g. the local extinction of a prey species).
These perturbations can afford an increase or decrease in population
size. For simplicity, we assume patches rapidly return to normal quality
after perturbations and we assume that the probability of death is
equal for all individuals so there is no change in the distribution
of individual states within patches. We assume these perturbations
have a normal distribution, although once again we must approximate
this normal distribution to interpolate transition probabilities on
to the grid of possible density states, where corresponding to the
range of possible population density transitions we define the set
$\boldsymbol{T}=\{-I,-I+1,...,I-1,I\}$ (where $I$ is the size of
the population density state grid), with a pseudo-normal distribution
over this set defining the probability of any given possible change
in density $\tau\in T$, defined as

\begin{equation}
\begin{array}{ccc}
p(\tau)=\frac{1}{K_{\gamma}}e^{-\frac{(\tau/I)^{2}}{2\sigma^{2}_{\gamma}}} & \text{for} & \tau\in T,\end{array}\label{eq:-2-1}
\end{equation}
where $\sigma^{2}_{\gamma}$ is the standard deviation parameter for
perturbations and $K_{\gamma}$ is 
\begin{equation}
K_{\gamma}=\underset{t=-I}{\overset{I}{\boldsymbol{\sum}}}e^{-\frac{(t/I)^{2}}{2\sigma^{2}_{\gamma}}}.\label{eq:-32}
\end{equation}
 The shift in probability generated by demographic feedback is defined
as 
\begin{equation}
c=round(\frac{q_{i}}{q_{I}}\cdot\gamma)\label{eq:-33}
\end{equation}
where $\gamma\in(0,1)$ is the maximum risk of mass mortality events
associated with population density. We can then define $p_{pert}(q_{i}|q_{j})$,
the probability of a patch of density j to a density of i as a result
of bounded patch quality perturbations, as
\begin{equation}
p_{pert}(q_{i}|q_{j})=\begin{cases}
\text{if }i=0 & \sum^{\tau=-I}_{\tau=min(I,i+c-j)}p(\tau)\\
\text{if }0<i<I & p(\tau=min(I,i+c-j)\\
\text{if }i=I & \sum^{\tau=I}_{\tau=min(I,i+c-j)}p(\tau)
\end{cases}\label{eq:-34}
\end{equation}

Given these perturbations, the proportion of patches that are of population
density $q_{i}$ for all $i$ in the next time step is 
\begin{equation}
h^{\prime\prime}(q_{i})=\sum^{I}_{j}p_{patch}(q_{i}|q_{j})\cdot h^{\prime}(q_{j})\label{eq:-35}
\end{equation}
and similarly the posterior physiological state distribution for each
patch state is $\forall i$
\[
\rho^{\prime\prime}_{i}(s,\pi,m)=\sum^{I}_{j}p_{patch}(q_{i}|q_{j})\cdot h^{\prime}(q_{j})\cdot\rho^{\prime\prime}_{i}(s,\pi,m).
\]


\subsection{Bayesian updating between growth cycles}

Individuals also know the transition probabilities between patch population
density states at demographic stability in their evolutionary history,
$\tilde{p}(q_{i}|q_{j})$, resulting from these growth and death processes,
and so individuals that survive to the next growth cycle will update
their information states accordingly for all $i$
\begin{equation}
\pi^{\prime}_{i}=\sum_{j}\pi_{j}\cdot\tilde{p}(q_{i}|q_{j}).\label{eq:-1}
\end{equation}
For the first instance, when we have not yet found demographic stability,
we will assume that their priors are that $\forall i\forall j,\tilde{p}(q_{i}|q_{j})=0$,
as this should not have a strong influence on the long term adaptive
dynamics assuming that the dynamics are ordered rather than chaotic.

\subsection{Demographic Stability}

Assuming that with successive generations $n$ that the distribution
of patches over states converges on stable values, i.e. $\forall q_{i}$
\begin{equation}
\text{\ensuremath{h^{*}(q_{i})}=lim}_{\rightarrow\infty}[h(q_{i})]\label{eq:-36}
\end{equation}
further, we can approximate $h^{*}(q_{i})$ by iterating over generations
until 
\begin{equation}
tol>\sum_{i}|h(q_{i})-h^{\prime\prime}(q_{i})|\label{eq:-37}
\end{equation}
where the tolerance threshold, $tol$, is very small, and $h^{\prime\prime}(q_{i})$
is the proportion of patches that are of population density state
$q_{i}$ after dispersal and deaths. Once demographic stability is
reached we assume $h^{*}(q_{i})=h(q_{i})$ and proceed with invasion
analysis.

\section{Finding Evolutionary Stability}

Having found the demographically stable state of the metapopulation,
given the resident strategy, we can then use the Bellman equation
to find the best response strategy $b(x(s,\pi,m))$ for an invading
individual. We then assume the best response strategy becomes the
new resident strategy, and recalculate the demographically stable
state of the metapopulation. The evolutionarily stable strategy (ESS)
is a resident strategy for which invading individuals can do no better
than to follow the resident strategy and so the ESS is defined as
\begin{equation}
\bar{x}(s,\pi,m)=b(x(s,\pi,m))=x(s,\pi,m).\label{eq:-38-1}
\end{equation}
Assuming there is a singular ESSs, we can iterating over invasions
till the best response strategy converges on the ESS, and the starting
point of the population will not matter. By sampling from multiple
starting points we can test this assumption, and if necessary plot
out the adaptive dynamics with phase portraits where there are multiple
ESSs or oscillatory dynamics.

To find the best response we define a Bellman fitness equation for
individuals in the model where $V_{n=0}(s,\pi,m)=0$ for all state
combinations (as the individual is dead), and iterate backwards to
the previous time step $n+1$
\begin{equation}
V_{n+1}(s,\pi,m)=R(s,m)+\text{max}_{x(s,\pi,m)=\kappa}[\sum_{s_{n}}\sum_{\pi_{n}}\sum_{m_{n}}a(s^{\prime},\pi^{\prime\prime},m^{\prime}|s,\pi,m,\kappa)\cdot V_{n}(s_{n},\pi_{n},m_{n})],\label{eq:-39-1}
\end{equation}
where $R(s,m)$ is the rate of reproduction for an individual of size
$s$ and quality $m$ (see equation 1), $p_{pert}(survival|\pi_{n})$
is the probability the individual survives between time-steps given
patch quality perturbations, and $p(s_{n},\pi_{n},m_{n},survives|s,\pi,m,\kappa)$
is the survival (given maintenance) and state transition probability
to the posterior state $\{s_{n},\pi_{n},m_{n}\}$ of individuals of
prior state $\{s,\pi,m\}$. 

\subsection{State transitions between developmental actions}

Individuals are born with accurate prior information regarding the
probability distribution over the population density states of patches
\begin{equation}
\boldsymbol{\pi}=\{\pi_{0}=p(q_{0})=h(q_{0}),\pi_{1}=p(q_{1})=h(q_{1}),...,\pi_{I}=p(q_{I})=h(q_{I})\}\label{eq:-3-1}
\end{equation}

Individuals have a probability of death from patch depletion and a
probability of death $\lambda$ from exogenous factors (e.g. old age).
Given an individual's information state after foraging $\boldsymbol{\pi}^{\prime}=\{\pi_{j=1},\pi_{j=2},...,\pi_{j=I}\}$
and given the probability of patch density transition at demographic
stability, the probability that an individual dies due to patch depletion
is 
\begin{equation}
P_{pert}(survival|\boldsymbol{\pi})=(1-\lambda)\cdot\sum_{j}\sum_{k}\pi_{j}\cdot(1-p_{pert}(q_{k}|q_{j})\cdot\frac{q_{j}-q_{k}}{q_{j}})\label{eq:-45}
\end{equation}
Individuals also know the transition probabilities between patch population
density states at demographic stability given patch perturbations
and dispersal, since the previous development action. The transition
probability at demographic stability of a patch of population density
$j$ to a density of $i$, given perturbations and dispersal is full
therefore, the posterior belief $\pi=\{\pi_{0},\pi_{1},\pi_{2}\}$becomes
\[
\pi^{\prime}_{0}=\pi_{0}\cdot\sum_{q_{j}}\tilde{p}_{births}(q_{0}|q_{j})\cdot\tilde{p}_{pert}(q_{j}|q_{0})+\pi_{1}\cdot\sum_{q_{j}}\tilde{p}_{births}(q_{0}|q_{j})\cdot\tilde{p}_{pert}(q_{j}|q_{1})+\pi_{2}\cdot\sum_{q_{j}}\tilde{p}_{births}(q_{0}|q_{j})\cdot\tilde{p}_{pert}(q_{j}|q_{2})
\]
\[
\pi^{\prime}_{1}=\pi_{0}\cdot\sum_{q_{j}}\tilde{p}_{births}(q_{1}|q_{j})\cdot\tilde{p}_{births}(q_{j}|q_{0})+\pi_{1}\cdot\sum_{q_{j}}\tilde{p}_{births}(q_{1}|q_{j})\cdot\tilde{p}_{births}(q_{j}|q_{1})+\pi_{2}\cdot\sum_{q_{j}}\tilde{p}_{births}(q_{1}|q_{j})\cdot\tilde{p}_{births}(q_{j}|q_{2})
\]
\[
\pi^{\prime}_{2}=\pi_{0}\cdot\sum_{q_{j}}\tilde{p}_{births}(q_{2}|q_{j})\cdot\tilde{p}_{births}(q_{j}|q_{0})+\pi_{1}\cdot\sum_{q_{j}}\tilde{p}_{births}(q_{2}|q_{j})\cdot\tilde{p}_{births}(q_{j}|q_{1})+\pi_{2}\cdot\sum_{q_{j}}\tilde{p}_{births}(q_{2}|q_{j})\cdot\tilde{p}_{births}(q_{j}|q_{2})
\]
ensuring that the posterior belief $\sum_{i}\pi^{\prime}_{i}=1$,
$\pi^{\prime}$ then has to be interpolated on the discrete belief
state space, generating posterior belief states $d_{1}$, $d_{2}$,
and $d_{3}$, with respective probabilities $p_{1}$, $p_{2}$, and
$p_{3}$. We can then generate the transition probability matrix for
belief states given that population growth and patch perturbations
will have occurred within the time between developmental actions,
i.e. $p(\pi^{\prime}|\pi,time)$, defined by running the following
algorithm for all prior belief states:
\begin{enumerate}
\item $\forall\pi^{\prime}$, $p(\pi^{\prime}|\pi,time)=0$
\item $\text{add \ensuremath{p_{1}}to }p(\pi^{\prime}=d_{1}|\pi,time)$,
\item $\text{add \ensuremath{p_{2}}to }p(\pi^{\prime}=d_{2}|\pi,time)$,
\item $\text{add \ensuremath{p_{3}}to }p(\pi^{\prime}=d_{3}|\pi,time)$\@.
\end{enumerate}
The posterior belief $\pi^{\prime}$ is then updated again upon receiving
a food mass of g for each of $i\in\{0,1,2\}$

\[
\pi^{\prime\prime}_{i}=\frac{\pi^{\prime}_{i}\cdot p(g|p(q_{i}))}{\sum^{I}_{k}\pi^{\prime}_{k}\cdot p(g_{k}|p(q_{k}))}
\]

ensuring that the posterior belief $\sum_{i}\pi^{\prime\prime}_{i}=1$,
$\pi^{\prime\prime}$ then has to be interpolated on the discrete
belief state space, generating posterior belief states $d_{1}$, $d_{2}$,
and $d_{3}$, with respective probabilities $p_{1}$, $p_{2}$, and
$p_{3}$. We can then generate the transition probability matrix for
belief states given a food mass of $g$ for each prior belief state
$\pi^{\prime}$
\begin{enumerate}
\item $\forall\pi^{\prime\prime}$, $p(\pi^{\prime\prime}|\pi^{\prime},g)=0$
\item $\text{add \ensuremath{p_{1}}to }p(\pi^{\prime\prime}=d_{1}|\pi^{\prime},g)$,
\item $\text{add \ensuremath{p_{2}}to }p(\pi^{\prime\prime}=d_{2}|\pi^{\prime},g)$,
\item $\text{add \ensuremath{p_{3}}to }p(\pi^{\prime\prime}=d_{3}|\pi^{\prime},g)$\@.
\end{enumerate}

\subsection{State Transitions During Development}

As observations of patch density and an individual's development are
linked by the same random variable, food mass $g$, we define the
state transition probability given individual decision making as
\begin{equation}
a(s^{\prime},\pi^{\prime\prime},m^{\prime}|s,\pi,m,\kappa)=P_{pert}(survival|\boldsymbol{\pi})\cdot\sum_{\pi^{\prime}}[p(\pi^{\prime}|\pi,time)\cdot\sum_{g}[p(g|\pi^{\prime})\cdot p(\pi^{\prime\prime}|\pi^{\prime},g)\cdot p(s^{\prime},m^{\prime},survives|\kappa,s,m,\pi^{\prime\prime},g)]]\label{eq:-40}
\end{equation}
where $p(s^{\prime},m^{\prime},survives|\kappa,s,m,\pi^{\prime\prime},g)$
is the probability of posterior size $s^{\prime}$ and quality $m^{\prime}$
given the individual's developmental policy dependent on their current
size, quality, and information states, given the food mass they received,
and given their probability of survival given maintenance costs, all
defined once by doing the following:
\begin{enumerate}
\item For all $\{s,\pi,m,g,\kappa\}$
\[
p(s^{\prime},m^{\prime}|s,\pi,m,g,\kappa)=0.
\]
\item For all $i$ and all $g$, define the transitions in state predicted
by their developmental strategy $\kappa,$where 
\begin{equation}
\Delta m=max(g-\kappa,0)\label{eq:-20-1-1}
\end{equation}
is the change in quality state of individuals after ingesting a food
mass of $g$, and where
\begin{equation}
\Delta s=max(g-\Delta m-y(s,m),0),\label{eq:-19-1-1}
\end{equation}
is the change in size state of individuals after ingesting a food
mass of $g$, given that all energy allocated to growth below the
fixed and state-dependent maintenance cost threshold $y(s,m)$, below
which all food mass is invested in maintenance.
\item Then add $p(survival|g,s,m,\kappa)$ to $p(s+\Delta s,m+\Delta m,survives|s,\pi,m,g,\kappa)$,
where depending on the food mass $g$ an individual encounters and
the fraction of that mass that they invest in maintenance, there is
a probability of death from lack of maintenance of 
\begin{equation}
P(survival|g,s,m,\kappa)=min(1,\frac{g-\Delta m-\Delta s}{y(s,m)})\label{eq:-44-1}
\end{equation}
where $y(s,m)$ is the size and quality dependent maintenance cost
threshold to avoid death completely, defined as shown previously:
\begin{equation}
y(s,m)=\frac{s^{3}\cdot m}{S^{3}\cdot M}\cdot g_{max}\label{eq:-22-1}
\end{equation}
\end{enumerate}

\end{document}
